\Needspace{19\baselineskip}

\CNsection{1}{Why Separate S-GW and P-GW?}

LTE deliberately splits the user-plane gateway into two distinct nodes:

{\def\LTcaptype{none} % do not increment counter
\Needspace{11\baselineskip}
\begin{longtable}[c]{@{}
  >{\raggedright\arraybackslash\hspace{0pt}}m{(\linewidth - 4\tabcolsep) * \real{0.3333}}
  >{\raggedright\arraybackslash\hspace{0pt}}m{(\linewidth - 4\tabcolsep) * \real{0.3333}}
  >{\raggedright\arraybackslash\hspace{0pt}}m{(\linewidth - 4\tabcolsep) * \real{0.3333}}@{}}
\toprule\noalign{}
\rowcolor{CNink}\CNth{Role} & \CNth{S-GW} & \CNth{P-GW} \\
\CNheadrulesep
\endhead
\bottomrule\noalign{}
\endlastfoot
\textbf{Anchoring} & Local mobility anchor (intra-LTE) & PDN/\allowbreak{}IP anchor (inter-RAT, global) \\
\textbf{Scope} & Per-MME pool area & Per-PDN connection (could span MME areas) \\
\textbf{Mobility} & Changes during inter-MME HO & Stays the same for entire PDN session \\
\textbf{Placement} & Close to RAN (low latency) & Centralized (policy, charging) \\
\end{longtable}
}

\textbf{Design Rationale:}

\begin{enumerate}
\def\labelenumi{\arabic{enumi}.}
\tightlist
\item
  \textbf{Different lifetimes}: S-GW can relocate during mobility; P-GW remains fixed for IP continuity
\item
  \textbf{Independent scaling}: S-GW scales with local traffic; P-GW scales with subscribers/\allowbreak{}policies
\item
  \textbf{CUPS preparation}: The split enables Control/\allowbreak{}User Plane Separation — you can distribute S-GW-U at the edge while keeping S-GW-C centralized
\item
  \textbf{Roaming}: In home-routed roaming, S-GW is in visited PLMN; P-GW is in home PLMN (S8 interface)
\end{enumerate}

\Needspace{14\baselineskip}

\CNsection{2}{S-GW (Serving Gateway)}

\CNchained\CNsubsection{2.1}{User-Plane Forwarding}

The S-GW is essentially a \textbf{GTP-U tunnel switch}:

\begin{itemize}
\tightlist
\item
  Receives DL packets from P-GW on S5/\allowbreak{}S8 GTP-U tunnel
\item
  Forwards them to eNodeB (eNB) on S1-U GTP-U tunnel
\item
  Receives UL packets from eNB on S1-U and forwards to P-GW on S5/\allowbreak{}S8
\end{itemize}

Each bearer has \textbf{separate GTP-U tunnels} on each interface (identified by TEID).

\CNsubsection{2.2}{Local Mobility Anchor}

During \textbf{X2 handover} (direct eNB-to-eNB):

\begin{itemize}
\tightlist
\item
  S-GW does NOT change
\item
  Source eNB forwards buffered/\allowbreak{}in-transit packets to target eNB via X2
\item
  After handover, MME sends \textbf{Modify Bearer Request} to S-GW with new eNB address/\allowbreak{}TEID
\item
  S-GW switches S1-U tunnel endpoint from source eNB to target eNB (\textbf{path switch})
\end{itemize}

This avoids tearing down the S5 tunnel to P-GW — only the S1-U leg changes.

\CNsubsection{2.3}{Downlink Data Buffering}

When UE is in \textbf{ECM-IDLE} (no S1-U tunnel exists):

\begin{enumerate}
\def\labelenumi{\arabic{enumi}.}
\tightlist
\item
  DL packets arrive from P-GW on S5 tunnel
\item
  S-GW has no S1-U endpoint \ensuremath{\to} \textbf{buffers packets}
\item
  S-GW sends \textbf{Downlink Data Notification} to MME (via S11)
\item
  MME initiates paging
\item
  After UE responds (Service Request), MME sends Modify Bearer Request with new eNB TEID
\item
  S-GW flushes buffered packets to eNB
\end{enumerate}

\textbf{Buffer limits}: Configurable per operator; typically 1-10 packets per bearer. Excess dropped.

\CNsubsection{2.4}{Per-UE/\allowbreak{}Per-Bearer Charging}

S-GW generates \textbf{CDRs (Charging Data Records)} with:

\begin{itemize}
\tightlist
\item
  UL/\allowbreak{}DL volume per bearer
\item
  Duration
\item
  QoS information (QCI, ARP, GBR/\allowbreak{}MBR)
\item
  Cause of record closure (HO, time limit, volume limit)
\end{itemize}

Sent to \textbf{offline charging system (CDF)} via Rf/\allowbreak{}Gz interface.

\CNsubsection{2.5}{Lawful Interception}

S-GW can mirror user-plane traffic for lawful interception (LI):

\begin{itemize}
\tightlist
\item
  X2 interface (ETSI standard) to Law Enforcement Monitoring Facility (LEMF)
\item
  Triggered by warrant, activated via X1 interface from LI administration
\item
  Duplicates packets without affecting UE service
\end{itemize}

\Needspace{17\baselineskip}

\CNsubsection{2.6}{S-GW Interfaces}

{\def\LTcaptype{none} % do not increment counter
\Needspace{13\baselineskip}
\begin{longtable}[c]{@{}
  >{\raggedright\arraybackslash\hspace{0pt}}m{(\linewidth - 6\tabcolsep) * \real{0.2973}}
  >{\raggedright\arraybackslash\hspace{0pt}}m{(\linewidth - 6\tabcolsep) * \real{0.1622}}
  >{\raggedright\arraybackslash\hspace{0pt}}m{(\linewidth - 6\tabcolsep) * \real{0.2703}}
  >{\raggedright\arraybackslash\hspace{0pt}}m{(\linewidth - 6\tabcolsep) * \real{0.2703}}@{}}
\toprule\noalign{}
\rowcolor{CNink}\CNth{Interface} & \CNth{Peer} & \CNth{Protocol} & \CNth{Function} \\
\CNheadrulesep
\endhead
\bottomrule\noalign{}
\endlastfoot
\textbf{S1-U} & eNB & GTP-U (UDP) & User-plane to/\allowbreak{}from RAN \\
\textbf{S5} & P-GW (same PLMN) & GTP-U + GTPv2-C & UP forwarding + CP session mgmt \\
\textbf{S8} & P-GW (roaming) & GTP-U + GTPv2-C & Same as S5 but inter-PLMN \\
\textbf{S11} & MME & GTPv2-C & CP: bearer/\allowbreak{}session management \\
\textbf{S12} & UTRAN (3G) & GTP-U & Direct tunnel for 3G interworking \\
\end{longtable}
}

\Needspace{14\baselineskip}

\CNsection{3}{P-GW (PDN Gateway)}

\CNchained\CNsubsection{3.1}{PDN Connectivity and IP Address Allocation}

The P-GW is the \textbf{UE’s gateway to the Internet/\allowbreak{}PDN} — it assigns the UE its IP address:

\textbf{IPv4 Allocation Methods:}

\begin{itemize}
\tightlist
\item
  \textbf{Static}: Pre-configured in HSS subscription data
\item
  \textbf{Dynamic (DHCPv4)}: P-GW acts as DHCP server or relays to external DHCP
\item
  \textbf{Dynamic (internal pool)}: P-GW allocates from local address pool (most common)
\end{itemize}

\textbf{IPv6 Allocation:}

\begin{itemize}
\tightlist
\item
  P-GW assigns a /\allowbreak{}64 prefix via \textbf{Router Advertisement (SLAAC)}
\item
  UE generates Interface ID (EUI-64 or privacy extensions)
\item
  Prefix delivered in PCO (Protocol Configuration Options) during bearer activation
\end{itemize}

\textbf{Dual-Stack (IPv4v6):}

\begin{itemize}
\tightlist
\item
  Both IPv4 address and IPv6 prefix assigned simultaneously
\item
  Requires PDN Type = IPv4v6 in subscription
\end{itemize}

\CNsubsection{3.2}{Policy Enforcement}

P-GW enforces policies received from \textbf{PCRF} via the \textbf{Gx interface}:

\begin{itemize}
\tightlist
\item
  \textbf{PCC Rules (Policy and Charging Control):}

  \begin{itemize}
  \tightlist
  \item
    SDF (Service Data Flow) filters — 5-tuple matching
  \item
    QoS parameters (QCI, GBR, MBR, ARP)
  \item
    Charging keys and rating groups
  \item
    Gate status (open/\allowbreak{}close specific flows)
  \end{itemize}
\item
  \textbf{TFT (Traffic Flow Template):}

  \begin{itemize}
  \tightlist
  \item
    \textbf{DL TFT}: At P-GW — maps incoming DL packets to correct bearer
  \item
    \textbf{UL TFT}: Sent to UE — UE maps UL packets to correct bearer
  \item
    Packet filters: source/\allowbreak{}dest IP, port range, protocol, ToS/\allowbreak{}DSCP
  \end{itemize}
\end{itemize}

\Needspace{12\baselineskip}

\CNsubsection{3.3}{Charging}

{\def\LTcaptype{none} % do not increment counter
\Needspace{8\baselineskip}
\begin{longtable}[c]{@{}
  >{\raggedright\arraybackslash\hspace{0pt}}m{(\linewidth - 6\tabcolsep) * \real{0.1818}}
  >{\raggedright\arraybackslash\hspace{0pt}}m{(\linewidth - 6\tabcolsep) * \real{0.3333}}
  >{\raggedright\arraybackslash\hspace{0pt}}m{(\linewidth - 6\tabcolsep) * \real{0.1818}}
  >{\raggedright\arraybackslash\hspace{0pt}}m{(\linewidth - 6\tabcolsep) * \real{0.3030}}@{}}
\toprule\noalign{}
\rowcolor{CNink}\CNth{Type} & \CNth{Interface} & \CNth{Node} & \CNth{Use Case} \\
\CNheadrulesep
\endhead
\bottomrule\noalign{}
\endlastfoot
\textbf{Online (OCS)} & Gy (Diameter) & OCS & Prepaid; real-time credit control \\
\textbf{Offline (OFCS)} & Gz (GTP’) & CDF/\allowbreak{}CGF & Postpaid; CDR generation \\
\end{longtable}
}

\textbf{Charging triggers:}

\begin{itemize}
\tightlist
\item
  Volume threshold (e.g., every 1~MB)
\item
  Time threshold (e.g., every 60 seconds)
\item
  QoS change, bearer modification
\item
  Re-authorization from OCS
\end{itemize}

\CNsubsection{3.4}{Packet Filtering and Routing}

\begin{itemize}
\tightlist
\item
  \textbf{DL}: Internet \ensuremath{\to} P-GW: Match destination IP to UE session, apply SDF filter to select bearer, encapsulate in GTP-U toward S-GW
\item
  \textbf{UL}: UE \ensuremath{\to} P-GW: Decapsulate GTP-U, apply NAT (if private addressing), route to Internet
\item
  \textbf{NAT/\allowbreak{}NAPT}: P-GW often performs NAT for IPv4 (private UE addresses \ensuremath{\to} public)
\item
  \textbf{Header compression}: Not at P-GW (done at PDCP layer in RAN)
\end{itemize}

\CNsubsection{3.5}{SGi Interface}

\begin{itemize}
\tightlist
\item
  SGi connects P-GW to \textbf{external PDNs} (Internet, IMS, enterprise VPNs)
\item
  IP-based interface (no GTP)
\item
  Equivalent to Gi in 3G (GGSN \ensuremath{\to} Internet)
\item
  Multiple PDNs supported (each identified by APN)
\item
  May connect to:

  \begin{itemize}
  \tightlist
  \item
    Internet (public APN)
  \item
    IMS core (ims APN)
  \item
    Enterprise VPN (private APN)
  \item
    Content/\allowbreak{}CDN platforms
  \end{itemize}
\end{itemize}

\CNkeepfig{m13-tft-mapping}{3}

\CNsubsection{3.6}{UL/\allowbreak{}DL TFT — Mapping Packets to Bearers}

\CNfigure{m13-tft-mapping}{Several bearers per PDN connection and how packets are mapped to them}

\Needspace{14\baselineskip}

\Needspace{31\baselineskip}

\CNsection{4}{GTP Protocol}

\CNchained\CNsubsection{4.1}{GTP-C (GTP Control Plane) — GTPv2-C}

GTP-C manages tunnel/\allowbreak{}session lifecycle on S11, S5/\allowbreak{}S8, S10, S3 interfaces.

\textbf{Key Messages:}

{\def\LTcaptype{none} % do not increment counter
\Needspace{18\baselineskip}
\begin{longtable}[c]{@{}
  >{\raggedright\arraybackslash\hspace{0pt}}m{(\linewidth - 4\tabcolsep) * \real{0.3103}}
  >{\raggedright\arraybackslash\hspace{0pt}}m{(\linewidth - 4\tabcolsep) * \real{0.3793}}
  >{\raggedright\arraybackslash\hspace{0pt}}m{(\linewidth - 4\tabcolsep) * \real{0.3103}}@{}}
\toprule\noalign{}
\rowcolor{CNink}\CNth{Message} & \CNth{Direction} & \CNth{Purpose} \\
\CNheadrulesep
\endhead
\bottomrule\noalign{}
\endlastfoot
Create Session Request/\allowbreak{}Response & MME\ensuremath{\to}S-GW\ensuremath{\to}P-GW & Establish PDN connection + default bearer \\
Modify Bearer Request/\allowbreak{}Response & MME\ensuremath{\to}S-GW & Update tunnel endpoints (HO, service request) \\
Delete Session Request/\allowbreak{}Response & MME\ensuremath{\to}S-GW\ensuremath{\to}P-GW & Tear down PDN connection \\
Create Bearer Request/\allowbreak{}Response & P-GW\ensuremath{\to}S-GW\ensuremath{\to}MME & Dedicated bearer activation \\
Update Bearer Request/\allowbreak{}Response & P-GW\ensuremath{\to}S-GW\ensuremath{\to}MME & QoS modification \\
Delete Bearer Request/\allowbreak{}Response & P-GW\ensuremath{\to}S-GW\ensuremath{\to}MME & Dedicated bearer deactivation \\
Downlink Data Notification & S-GW\ensuremath{\to}MME & Trigger paging for idle UE \\
Release Access Bearers & MME\ensuremath{\to}S-GW & UE goes idle, release S1-U \\
\end{longtable}
}

\CNkeepfig{m13-gtpv2c-header}{2}

\textbf{GTPv2-C Header:}

\CNfigure{m13-gtpv2c-header}{GTPv2-C header}

\textbf{TEID (Tunnel Endpoint Identifier):}

\begin{itemize}
\tightlist
\item
  32-bit identifier, locally unique per node
\item
  Each side assigns its own TEID and communicates it to the peer
\item
  Allows multiplexing multiple tunnels on same IP:port
\end{itemize}

\CNkeepfig{m13-gtpu-header}{7}

\CNsubsection{4.2}{GTP-U (GTP User Plane)}

GTP-U encapsulates user IP packets for transport between EPC nodes.

\textbf{GTP-U Header:}

\CNfigure{m13-gtpu-header}{GTP-U header}

\begin{itemize}
\tightlist
\item
  \textbf{Message Type 0xFF (255)}: T-PDU — contains user data (most common)
\item
  \textbf{Transport}: UDP port 2152
\item
  \textbf{T-PDU}: The encapsulated user IP packet
\item
  \textbf{Sequence Number}: Optional, used for reordering during path switch/\allowbreak{}handover
\end{itemize}

\Needspace{25\baselineskip}

\CNsubsection{4.3}{GTP-C vs GTP-U Comparison}

{\def\LTcaptype{none} % do not increment counter
\Needspace{21\baselineskip}
\begin{longtable}[c]{@{}
  >{\raggedright\arraybackslash\hspace{0pt}}m{(\linewidth - 4\tabcolsep) * \real{0.1860}}
  >{\raggedright\arraybackslash\hspace{0pt}}m{(\linewidth - 4\tabcolsep) * \real{0.3953}}
  >{\raggedright\arraybackslash\hspace{0pt}}m{(\linewidth - 4\tabcolsep) * \real{0.4186}}@{}}
\toprule\noalign{}
\rowcolor{CNink}\CNth{Aspect} & \CNth{GTP-C (GTPv2-C)} & \CNth{GTP-U (GTPv1-U)} \\
\CNheadrulesep
\endhead
\bottomrule\noalign{}
\endlastfoot
\textbf{Purpose} & Session/\allowbreak{}tunnel management & User data transport \\
\textbf{Version} & v2 (in LTE) & v1 (unchanged from 3G) \\
\textbf{Transport} & UDP port 2123 & UDP port 2152 \\
\textbf{Reliability} & Request/\allowbreak{}Response + retransmission & No ACK (relies on upper layers) \\
\textbf{TEID} & Identifies control session & Identifies user data tunnel (per-bearer) \\
\textbf{Messages} & Create/\allowbreak{}Modify/\allowbreak{}Delete Session, etc. & T-PDU (0xFF), Echo, Error Indication \\
\textbf{Interfaces} & S11, S5/\allowbreak{}S8, S10, S3 & S1-U, S5/\allowbreak{}S8 (user plane) \\
\textbf{Payload} & IE (Information Elements) structured & Raw IP packet (T-PDU) \\
\textbf{Header size} & 12~bytes (min) & 8~bytes (min), 12 with options \\
\textbf{Statefulness} & Stateful (sessions) & Stateless (per-packet forwarding) \\
\end{longtable}
}

\Needspace{14\baselineskip}

\CNkeepfig{m13-downlink-trace}{8}

\CNsection{5}{Packet Path Trace (End-to-End with Encapsulation)}

\CNchained\CNsubsection{5.1}{Downlink: Internet \ensuremath{\to} UE}

\CNfigure{m13-downlink-trace}{Downlink packet trace from the Internet to the UE with each encapsulation}

\CNkeepfig{m13-headers}{3}

\CNsubsection{5.2}{Headers at Each Point}

\CNfigure{m13-headers}{Headers of a downlink packet on each segment}

\CNkeepfig{m13-gtp-tunnels}{3}

\CNsubsection{5.3}{GTP Tunnel Structure Diagram}

\CNfigure{m13-gtp-tunnels}{GTP tunnel endpoints and TEIDs along the user-plane path}

\textbf{Note:} Each GTP-U tunnel segment has \textbf{two TEIDs} — one for each direction:

\begin{itemize}
\tightlist
\item
  S1-U UL: eNB\ensuremath{\to}S-GW uses S-GW’s TEID (0x5678)
\item
  S1-U DL: S-GW\ensuremath{\to}eNB uses eNB’s TEID (0x1234)
\end{itemize}

\Needspace{14\baselineskip}

\CNkeepfig{m13-cups-split}{10}

\CNsection{6}{CUPS: Control and User Plane Separation}

\CNchained\CNsubsection{6.1}{Concept (3GPP Release 14+)}

CUPS splits S-GW and P-GW each into control-plane and user-plane components:

\CNfigure{m13-cups-split}{CUPS splits S-GW and P-GW into control-plane and user-plane parts}

\CNsubsection{6.2}{Sx Interface and PFCP Protocol}

\textbf{PFCP (Packet Forwarding Control Protocol)} — 3GPP TS 29.244:

\begin{itemize}
\tightlist
\item
  Runs between CP and UP components (Sxa for S-GW, Sxb for P-GW, Sxc for TDF)
\item
  \textbf{Session Establishment}: CP tells UP to create forwarding rules
\item
  \textbf{Session Modification}: CP updates rules (e.g., new TEID after handover)
\item
  \textbf{Session Deletion}: CP instructs UP to remove forwarding state
\end{itemize}

\textbf{PFCP Key Concepts:}\CNbr{}\textbar{} Element \textbar{} Purpose \textbar{}\CNbr{}\textbar{}———\textbar{}———\textbar{}\CNbr{}\textbar{} PDR (Packet Detection Rule) \textbar{} Match criteria (TEID, IP, port) \textbar{}\CNbr{}\textbar{} FAR (Forwarding Action Rule) \textbar{} What to do (forward, buffer, drop) \textbar{}\CNbr{}\textbar{} QER (QoS Enforcement Rule) \textbar{} Rate limiting (GBR/\allowbreak{}MBR) \textbar{}\CNbr{}\textbar{} URR (Usage Reporting Rule) \textbar{} Volume/\allowbreak{}time measurement for charging \textbar{}\CNbr{}\textbar{} BAR (Buffering Action Rule) \textbar{} Buffering behavior for idle UE \textbar{}

\CNsubsection{6.3}{Benefits of CUPS}

\begin{enumerate}
\def\labelenumi{\arabic{enumi}.}
\tightlist
\item
  \textbf{Distributed user plane}: Deploy S-GW-U /\allowbreak{} P-GW-U at edge (low latency for MEC)
\item
  \textbf{Centralized control}: Keep S-GW-C /\allowbreak{} P-GW-C in central DC (simpler management)
\item
  \textbf{Independent scaling}: Scale UP with traffic volume, CP with signaling load
\item
  \textbf{Flexible deployment}: Multiple UP instances per CP, or vice versa
\item
  \textbf{5G readiness}: CUPS architecture directly maps to 5G SMF (CP) + UPF (UP) split
\end{enumerate}

\CNkeepfig{m13-cups}{3}

\CNsubsection{6.4}{CUPS Architecture Diagram}

\CNfigure{m13-cups}{CUPS: S-GW and P-GW split into control and user planes}

\textbf{Blue = Control Plane} \textbar{} \textbf{Green = User Plane}

\Needspace{14\baselineskip}

\Needspace{27\baselineskip}

\CNsection{7}{Evolution: S-GW + P-GW \ensuremath{\to} UPF (5G)}

\CNchained\CNsubsection{7.1}{Why Merged into UPF?}

In 5G, S-GW and P-GW user-plane functions are \textbf{merged into a single UPF (User Plane Function)}:

{\def\LTcaptype{none} % do not increment counter
\Needspace{16\baselineskip}
\begin{longtable}[c]{@{}
  >{\raggedright\arraybackslash\hspace{0pt}}m{(\linewidth - 4\tabcolsep) * \real{0.2105}}
  >{\raggedright\arraybackslash\hspace{0pt}}m{(\linewidth - 4\tabcolsep) * \real{0.5000}}
  >{\raggedright\arraybackslash\hspace{0pt}}m{(\linewidth - 4\tabcolsep) * \real{0.2895}}@{}}
\toprule\noalign{}
\rowcolor{CNink}\CNth{Aspect} & \CNth{LTE (S-GW + P-GW)} & \CNth{5G (UPF)} \\
\CNheadrulesep
\endhead
\bottomrule\noalign{}
\endlastfoot
Nodes & 2 separate (S-GW, P-GW) & 1 unified (UPF) \\
GTP hops & 2 (S1-U + S5) & 1 (N3 only, or N3+N9 for chaining) \\
Control & S-GW-C + P-GW-C & SMF (single CP for UP) \\
CP-UP protocol & Sxa/\allowbreak{}Sxb PFCP & N4 PFCP (same protocol, evolved) \\
Deployment & Typically centralized & Distributed (edge UPF, anchor UPF) \\
Mobility anchor & Fixed P-GW & Flexible (SSC modes 1/\allowbreak{}2/\allowbreak{}3) \\
Slicing & N/\allowbreak{}A & UPF per slice \\
\end{longtable}
}

\CNsubsection{7.2}{Why the Merge Makes Sense}

\begin{enumerate}
\def\labelenumi{\arabic{enumi}.}
\tightlist
\item
  \textbf{Eliminates one GTP-U hop}: S1-U + S5 \ensuremath{\to} just N3 (+ optional N9 between UPFs)

  \begin{itemize}
  \tightlist
  \item
    Lower latency, fewer encap/\allowbreak{}decap operations
  \end{itemize}
\item
  \textbf{Simplified forwarding}: One node handles both local mobility and PDN anchoring
\item
  \textbf{Edge deployment}: UPF placed close to user for MEC/\allowbreak{}URLLC — no need for separate S-GW-U at edge
\item
  \textbf{Flexible anchoring (SSC modes)}:

  \begin{itemize}
  \tightlist
  \item
    SSC 1: Fixed UPF (like LTE P-GW) — IP continuity
  \item
    SSC 2: UPF can change — new IP address (for non-persistent sessions)
  \item
    SSC 3: Multi-homing — old + new UPF active simultaneously during transition
  \end{itemize}
\item
  \textbf{UPF chaining}: Multiple UPFs in series (I-UPF as intermediate, PSA-UPF as anchor) replaces S-GW\ensuremath{\to}P-GW chain with more flexible topology
\end{enumerate}

\Needspace{17\baselineskip}

\CNsubsection{7.3}{Interface Mapping}

{\def\LTcaptype{none} % do not increment counter
\Needspace{13\baselineskip}
\begin{longtable}[c]{@{}cc@{}}
\toprule\noalign{}
\rowcolor{CNink}\CNth{LTE} & \CNth{5G} \\
\CNheadrulesep
\endhead
\bottomrule\noalign{}
\endlastfoot
S1-U (eNB \ensuremath{\leftrightarrow} S-GW) & N3 (gNB \ensuremath{\leftrightarrow} UPF) \\
S5/\allowbreak{}S8 (S-GW \ensuremath{\leftrightarrow} P-GW) & N9 (UPF \ensuremath{\leftrightarrow} UPF, if chained) \\
SGi (P-GW \ensuremath{\leftrightarrow} Internet) & N6 (UPF \ensuremath{\leftrightarrow} DN) \\
S11 (MME \ensuremath{\leftrightarrow} S-GW) & N4 (SMF \ensuremath{\leftrightarrow} UPF) [PFCP] \\
Gx (PCRF \ensuremath{\leftrightarrow} P-GW) & N7 (PCF \ensuremath{\leftrightarrow} SMF) [SBI] \\
\end{longtable}
}

\Needspace{28\baselineskip}

\CNsection{}{Summary}

{\def\LTcaptype{none} % do not increment counter
\Needspace{22\baselineskip}
\begin{longtable}[c]{@{}ccc@{}}
\toprule\noalign{}
\rowcolor{CNink}\CNth{Function} & \CNth{S-GW} & \CNth{P-GW} \\
\CNheadrulesep
\endhead
\bottomrule\noalign{}
\endlastfoot
User-plane forwarding & \CNyes{} & \CNyes{} \\
GTP-U termination & S1-U + S5 & S5 + SGi \\
IP address allocation & \CNno{} & \CNyes{} \\
Policy enforcement (PCC) & \CNno{} & \CNyes{} \\
Charging (online) & \CNno{} & \CNyes{} (Gy) \\
Charging (offline) & \CNyes{} (CDRs) & \CNyes{} (Gz) \\
DL buffering + paging trigger & \CNyes{} & \CNno{} \\
Local mobility anchor & \CNyes{} & \CNno{} \\
PDN/\allowbreak{}IP anchor & \CNno{} & \CNyes{} \\
Lawful interception & \CNyes{} & \CNyes{} \\
NAT & \CNno{} & \CNyes{} \\
TFT/\allowbreak{}SDF enforcement & \CNno{} & \CNyes{} \\
\end{longtable}
}
